\documentclass[a4paper,twoside]{article}

\usepackage{morewrites}

\usepackage[svgnames,dvipsnames,x11names]{xcolor}
\usepackage{graphicx}
\usepackage[english]{babel}
\usepackage[colorlinks=true, linkcolor=NavyBlue, urlcolor=NavyBlue, citecolor=NavyBlue]{hyperref}
\usepackage[a4paper]{geometry}
\usepackage{ts-fonts}
\usepackage{amsmath}
\usepackage{float}
\usepackage{seqsplit}
\usepackage{ts-callouts}
\usepackage{ts-code}

\usepackage{lastpage}
\usepackage{refcount}
\makeatletter
\def\ps@plain{
  \let\@oddhead\@empty
  \let\@evenhead\@empty
  \def\@oddfoot{
    \hfil
    \ifnum\getpagerefnumber{LastPage}>1\relax
      \thepage
    \fi
    \hfil}
  \let\@evenfoot\@oddfoot
}
\makeatother

\title{Geometry Fragment}
\author{}\date{}

\begin{document}

\maketitle

\thispagestyle{plain}
The geometry fragment allows you to customize the page layout of your document, including paper size, orientation, margins, and adding watermarks. It relies on the \LaTeX{} \texttt{geometry} package to manage these settings and uses TikZ for watermarking if you want any.

\section{Paper format}\label{paper-format}

We support standard paper formats recognized by the \LaTeX{} \texttt{geometry} package. You can specify the format using the \texttt{paper.format} key. For example, to set the paper size to A5, you would use:

\begin{code}{yaml}{}{}
\begin{Verbatim}[commandchars=\\\{\},breaklines, breakanywhere, commandchars=\\\{\}]
\PY{n+nt}{press}\PY{p}{:}
\PY{+w}{  }\PY{n+nt}{paper}\PY{p}{:}\PY{+w}{ }\PY{l+lScalar+lScalarPlain}{a5}
\end{Verbatim}
\end{code}
The supported formats include but are not limited to: a0, a1, a2, a3, a4, a5, a6, b0, b1, b2, b3, b4, b5, b6, c0, c1, c2, c3, c4, c5, c6, letter, legal, executive, ansia, ansib, ansic, ansid, ansie.

\begin{callout}[callout note]{Note}
The default format is \texttt{a4}, in contrast to \LaTeX{}'s \texttt{letter} default.
Globally, only the United States, Canada, Mexico, and a few Caribbean
countries primarily use the \texttt{letter} size -- roughly 500 million people. The
rest of the world, representing more than 6 billion people, relies on \texttt{a4}
as the standard paper size. Given this overwhelming majority, TeXSmith
defaults to \texttt{a4} to better serve its global user base. Sorry, folks in the
US, Canada, and Mexico -- TeXSmith is opinionated and has chosen the
broadest consensus!
\end{callout}
\section{Orientation}\label{orientation}

You can set the page orientation using the \texttt{paper.orientation} key. The possible values are \texttt{portrait} and \texttt{landscape}. For example, to set the orientation to landscape, you would use:

\begin{code}{yaml}{}{}
\begin{Verbatim}[commandchars=\\\{\},breaklines, breakanywhere, commandchars=\\\{\}]
\PY{n+nt}{press}\PY{p}{:}
\PY{+w}{  }\PY{n+nt}{paper}\PY{p}{:}
\PY{+w}{    }\PY{n+nt}{orientation}\PY{p}{:}\PY{+w}{ }\PY{l+lScalar+lScalarPlain}{landscape}
\end{Verbatim}
\end{code}
\section{Margins}\label{margins}

You can customize the page margins using the \texttt{paper.margin} key. You can either specify a single value for all margins or provide an object with specific margins. For example:

\begin{code}{yaml}{}{}
\begin{Verbatim}[commandchars=\\\{\},breaklines, breakanywhere, commandchars=\\\{\}]
\PY{n+nt}{press}\PY{p}{:}
\PY{+w}{  }\PY{n+nt}{paper}\PY{p}{:}
\PY{+w}{    }\PY{n+nt}{margin}\PY{p}{:}\PY{+w}{ }\PY{l+lScalar+lScalarPlain}{2cm}
\end{Verbatim}
\end{code}
Or for specific margins:

\begin{code}{yaml}{}{}
\begin{Verbatim}[commandchars=\\\{\},breaklines, breakanywhere, commandchars=\\\{\}]
\PY{n+nt}{press}\PY{p}{:}
\PY{+w}{  }\PY{n+nt}{paper}\PY{p}{:}
\PY{+w}{    }\PY{n+nt}{margin}\PY{p}{:}
\PY{+w}{      }\PY{n+nt}{left}\PY{p}{:}\PY{+w}{ }\PY{l+lScalar+lScalarPlain}{3cm}
\PY{+w}{      }\PY{n+nt}{right}\PY{p}{:}\PY{+w}{ }\PY{l+lScalar+lScalarPlain}{2cm}
\PY{+w}{      }\PY{n+nt}{top}\PY{p}{:}\PY{+w}{ }\PY{l+lScalar+lScalarPlain}{4cm}
\PY{+w}{      }\PY{n+nt}{bottom}\PY{p}{:}\PY{+w}{ }\PY{l+lScalar+lScalarPlain}{5cm}
\end{Verbatim}
\end{code}
Or use predefined margin settings like \texttt{narrow}, \texttt{moderate}, or \texttt{wide}:

\begin{code}{yaml}{}{}
\begin{Verbatim}[commandchars=\\\{\},breaklines, breakanywhere, commandchars=\\\{\}]
\PY{n+nt}{press}\PY{p}{:}
\PY{+w}{  }\PY{n+nt}{paper}\PY{p}{:}
\PY{+w}{    }\PY{n+nt}{margin}\PY{p}{:}\PY{+w}{ }\PY{l+lScalar+lScalarPlain}{narrow}
\end{Verbatim}
\end{code}
\section{Examples}\label{examples}

Here is an example configuration that sets the paper size to C5, uses landscape orientation, adds a frame, customizes the margins, and includes a watermark:

\begin{code}{md}{}{}
\begin{Verbatim}[commandchars=\\\{\},breaklines, breakanywhere, commandchars=\\\{\}]
\PYZhy{}\PYZhy{}\PYZhy{}
press:
  paper:
    format: c5
    orientation: landscape
    frame: true
    margin:
      left: 3cm
      bottom: 5
\PY{g+gu}{    watermark: \PYZdq{}ENVELOPE\PYZdq{}}
\PY{g+gu}{\PYZhy{}\PYZhy{}\PYZhy{}}
\PY{g+gh}{\PYZsh{} Custom Geometry Example}

This document demonstrates custom page geometry settings using the geometry fragment.
\end{Verbatim}
\end{code}
\begin{figure}[H]
    \centering

    \ifdefined\adjustbox

        \adjustbox{max width=\textwidth}{
            \includegraphics[width=0.7\linewidth]{snippet-<HASH>.pdf}
        }

    \else

        \includegraphics[width=0.7\linewidth]{snippet-<HASH>.pdf}

    \fi

\end{figure}
Here's another example with custom paper width:

\begin{code}{md}{}{}
\begin{Verbatim}[commandchars=\\\{\},breaklines, breakanywhere, commandchars=\\\{\}]
\PYZhy{}\PYZhy{}\PYZhy{}
press:
  paper:
    width: 12cm
    height: 5cm
    frame: true
\PY{g+gu}{    margin: narrow}
\PY{g+gu}{\PYZhy{}\PYZhy{}\PYZhy{}}
\PYZdl{}\PYZdl{} E=mc\PYZca{}2 \PYZdl{}\PYZdl{}
\end{Verbatim}
\end{code}
\begin{figure}[H]
    \centering

    \ifdefined\adjustbox

        \adjustbox{max width=\textwidth}{
            \includegraphics[width=0.7\linewidth]{snippet-<HASH>.pdf}
        }

    \else

        \includegraphics[width=0.7\linewidth]{snippet-<HASH>.pdf}

    \fi

\end{figure}

\end{document}
